\documentclass[11pt,a4paper]{article}

% ---------------------------------------------------------
% PREAMBLE (stable, warning-free, Overleaf-safe)
% ---------------------------------------------------------
\usepackage[utf8]{inputenc}
\usepackage[T1]{fontenc}
\usepackage{lmodern}
\usepackage{amsmath, amssymb, amsfonts}
\usepackage{bm}
\usepackage{geometry}
\usepackage{graphicx}
\usepackage{hyperref}
\usepackage{cite}
\usepackage{authblk}
\usepackage{physics}
\usepackage{mathtools}

\geometry{margin=2.4cm}

\title{\textbf{Companion LXVIII: Cross-Correlation Functions of Persistence Diagrams in the Relaxation-Driven Cyclic Cosmology}}
\author{Michael Lehmann \\ \texttt{mi.lehmann@gmx.de}}
\date{April 2026}

\begin{document}
\maketitle

\begin{abstract}
Cross-correlation functions of persistence diagrams quantify the coupling of 
topological features across different filtration parameters, smoothing scales, or 
redshift slices. In weak-lensing analyses, these correlations measure how the 
birth and death of connected components, loops, and void-like structures co-evolve 
across the gravitational field. In the standard cosmological model, cross-
correlation functions reflect the nonlinear matter distribution and the projection 
of large-scale structure. In the Relaxation-Driven Cyclic Cosmology (RDCC), the 
scale-dependent evolution of the gravitational potential modifies the filtration-
coupled topology in a characteristic way. This Companion develops a continuous 
description of cross-correlation functions in RDCC, deriving how the relaxation 
scale influences topological coupling, persistence-lifetime alignment, and the 
observational signatures in weak-lensing surveys. The resulting framework provides 
a distinctive probe of the RDCC gravitational field through filtration-coupled 
topology.
\end{abstract}
% ---------------------------------------------------------
\section{Introduction}

Persistence diagrams encode the birth and death of topological features in 
weak-lensing convergence fields. While cross-covariance measures the correlated 
spread of topological features across ensembles, cross-correlation functions 
quantify the coupling of topological structure across different filtration 
parameters, smoothing scales, or redshift slices.

In the standard cosmological model, cross-correlation functions reflect the 
nonlinear matter distribution and the projection of large-scale structure. In 
RDCC, the gravitational potential evolves differently. The relaxation mechanism 
suppresses the decay of small-scale modes and enhances the coherence of large-scale 
modes. This modifies the filtration-coupled topology in a scale-dependent manner, 
producing a distinctive signature in cross-correlation analyses.
% ---------------------------------------------------------
\section{Cross-Correlation Functions of Persistence Diagrams}

Given two filtrations with persistence diagrams $D(\nu_{1})$ and $D(\nu_{2})$, the 
cross-correlation function is defined as

\begin{equation}
    \xi_{D}(\nu_{1}, \nu_{2})
    = \frac{
        \left\langle
            W_{p}(D(\nu_{1}), D^{*})
            W_{p}(D(\nu_{2}), D^{*})
        \right\rangle
      }{
        \sqrt{
            \left\langle W_{p}(D(\nu_{1}), D^{*})^{2} \right\rangle
            \left\langle W_{p}(D(\nu_{2}), D^{*})^{2} \right\rangle
        }
      }.
\end{equation}

This function measures the degree of filtration-coupled topology.
% ---------------------------------------------------------
\section{RDCC Modification of Persistence Diagrams}

In RDCC, the convergence evolves as
\begin{equation}
    \kappa_{\mathrm{RDCC}}(\ell)
    = \kappa(\ell)
      \left[
        1 - \chi_{\kappa}
        \left(\frac{\ell}{\ell_{\mathrm{rel}}}\right)^{\alpha}
      \right].
\end{equation}

This modifies the birth and death thresholds of topological features:

\begin{align}
    b_{\mathrm{RDCC}} &= b \left[ 1 - \chi_{b} (\ell/\ell_{\mathrm{rel}})^{\alpha} \right], \\
    d_{\mathrm{RDCC}} &= d \left[ 1 - \chi_{d} (\ell/\ell_{\mathrm{rel}})^{\alpha} \right].
\end{align}

Thus the RDCC filtration diagrams become
\begin{equation}
    D_{\mathrm{RDCC}}(\nu)
    = \left\{
        (b_{\mathrm{RDCC}}(\nu), d_{\mathrm{RDCC}}(\nu))
      \right\}.
\end{equation}
% ---------------------------------------------------------
\section{Cross-Correlation Functions in RDCC}

The RDCC cross-correlation function is

\begin{equation}
    \xi_{D,\mathrm{RDCC}}(\nu_{1}, \nu_{2})
    = \xi_{D}(\nu_{1}, \nu_{2})
      \left[
        1 - \chi_{\xi}
        \left(\frac{\ell}{\ell_{\mathrm{rel}}}\right)^{\alpha}
      \right].
\end{equation}

This modification reflects the relaxation-driven change in filtration-coupled 
topology.
% ---------------------------------------------------------
\section{Filtration-Coupled Topological Signatures of RDCC}

The relaxation mechanism enhances the coherence of large-scale modes. Cross-
correlation functions therefore exhibit:

\begin{itemize}
    \item enhanced correlation between large-scale filtration thresholds,
    \item suppressed correlation between small-scale thresholds,
    \item a characteristic turnover near $\ell_{\mathrm{rel}}$,
    \item modified redshift evolution of filtration-coupled topology,
    \item altered alignment of persistence lifetimes across thresholds.
\end{itemize}

These signatures provide a direct probe of the RDCC gravitational field through 
filtration-coupled topology.
% ---------------------------------------------------------
\section{Conclusion}

Cross-correlation functions of persistence diagrams in the Relaxation-Driven Cyclic 
Cosmology reflect the scale-dependent evolution of the gravitational potential and 
the nonlinear matter distribution. The relaxation mechanism modifies the 
filtration-coupled topology in a scale-dependent manner, producing distinctive 
signatures in weak-lensing surveys. These effects provide a direct observational 
window into the relaxation scale and the temporal structure of the RDCC 
gravitational field. The present Companion establishes the theoretical foundation 
for future studies of filtration-coupled topology, nonlinear structure, and the 
interplay between light and gravity in RDCC.

\bibliographystyle{apsrev4-2}
\nocite{*}
\bibliography{references}


\end{document}
